%% Question bank chapter file - generated by bank_tools.py build (do not edit while the bank is still being built:
%% it is regenerated from the per-test data in Question Bank\data). Every question is one bankq block:
%%   \begin{bankq}[id=..., source={...}, paper=.., qno=.., type=..., subtopic={...}, answer={...}, answer_src=key|solved,
%%                 working={...}, difficulty=1-10, diff_parts={K S I C F}, diff_note={...},
%%                 group=<passage id or empty>, also_in={duplicate copies in other tests}]
%% difficulty = min(10, K+S+I+C+F-1): K concepts 1-3, S steps 1-4, I insight 0-2, C computation 0-2, F format 0-1
%%   ...question LaTeX, options, figures (\bfig / \bfigi -> img/)...
%%   \end{bankq}
%% Passages shared by several questions are bankpassage blocks; their questions carry group=<passage id>.
\documentclass[10pt]{article}
\usepackage[a4paper,margin=1.8cm]{geometry}
\usepackage[T1]{fontenc}
\usepackage{amsmath,amssymb,mathtools}
\usepackage{graphicx}
\usepackage[export]{adjustbox}
\usepackage{xcolor,needspace,keyval}
\usepackage{enumitem}
\usepackage{booktabs,array,tabularx}
\usepackage[hidelinks]{hyperref}
\setlength{\parindent}{0pt}
\setlength{\parskip}{3pt}
\makeatletter
\define@key{bq}{id}{\def\bq@id{#1}}\define@key{bq}{source}{\def\bq@src{#1}}\define@key{bq}{paper}{\def\bq@paper{#1}}
\define@key{bq}{qno}{\def\bq@qno{#1}}\define@key{bq}{type}{\def\bq@type{#1}}\define@key{bq}{subtopic}{\def\bq@sub{#1}}
\define@key{bq}{answer}{\def\bq@ans{#1}}\define@key{bq}{answer_src}{\def\bq@asrc{#1}}\define@key{bq}{working}{\def\bq@work{#1}}
\define@key{bq}{group}{\def\bq@group{#1}}\define@key{bq}{also_in}{\def\bq@also{#1}}
\define@key{bq}{difficulty}{\def\bq@diff{#1}}\define@key{bq}{diff_parts}{\def\bq@dparts{#1}}\define@key{bq}{diff_note}{\def\bq@dnote{#1}}
\newenvironment{bankq}[1][]{\def\bq@work{}\def\bq@also{}\def\bq@group{}\def\bq@diff{?}\def\bq@dparts{}\def\bq@dnote{}\setkeys{bq}{#1}%
  \par\needspace{4\baselineskip}\medskip\noindent{\small\textbf{\bq@id}\hfill\color{gray}[\bq@src, Paper \bq@paper, Q\bq@qno]
  \bq@type\ $\cdot$ \bq@sub\ $\cdot$ \textbf{D\bq@diff}}\par\nopagebreak\relax}
  {\par{\small\color{teal}Answer: \textbf{\bq@ans}\ifx\bq@work\@empty\else\ (solved: \bq@work)\fi
  \ifx\bq@also\@empty\else\par\color{gray}Also in: \bq@also\fi}\par}
\define@key{bp}{id}{\def\bp@id{#1}}\define@key{bp}{source}{\def\bp@src{#1}}\define@key{bp}{questions}{\def\bp@qs{#1}}
\newenvironment{bankpassage}[1][]{\setkeys{bp}{#1}\par\needspace{4\baselineskip}\bigskip\noindent
  {\small\textbf{Passage \bp@id}\hfill\color{gray}[\bp@src; questions \bp@qs]}\par\nopagebreak\relax\itshape}{\par}
\makeatother
\newcommand{\bfig}[2]{\par\smallskip{\centering\includegraphics[width=#1pt,max width=0.62\linewidth,max height=0.32\textheight]{img/#2}\par}\smallskip}
\newcommand{\bfigi}[2]{\includegraphics[width=#1pt,max width=\linewidth,max height=0.32\textheight,valign=c]{img/#2}}
\newcommand{\bmissing}[1]{\textit{[figure missing in the source conversion: \detokenize{#1}]}}
\newenvironment{qoptions}{\begin{itemize}[leftmargin=1.8em,itemsep=1pt,topsep=2pt,label={}]}{\end{itemize}}
\newcommand{\opt}[2]{\item[(#1)] #2}
\providecommand{\fwstart}{}\providecommand{\fwend}{}\providecommand{\qfullsplit}{}
\newcolumntype{L}[1]{>{\raggedright\arraybackslash}p{#1}}
\newcolumntype{C}[1]{>{\centering\arraybackslash}p{#1}}
\newcolumntype{R}[1]{>{\raggedleft\arraybackslash}p{#1}}
\providecommand{\sectionhdr}[1]{}\providecommand{\subsectionhdr}[1]{}


\begin{document}
{\Large\bfseries Semiconductors --- Question Bank}\par\smallskip
{\small Physics chapter P24. Questions from all converted test papers, grouped by source series. Answers: printed key where it matches the paper, otherwise solved (answer\_src=solved, with working). Exact repeats across papers appear once, with also\_in listing the other copies.}\par\medskip\hrule

\section*{Score-II papers}
\begin{bankq}[id=SII05-P1-PHY-S1-Q11, source={Score-II - Test 05 (22 Mar 2022)}, paper=1, qno=11, type=single, subtopic={Logic gates}, answer={C}, answer_src=key, difficulty=4, diff_parts={K2 S1 I1 C1 F0}, diff_note={concepts: standard problem in this chapter, logic gates | steps: apply the relevant results; compute}]
Find output Y :
\bfig{143.8}{SII05_p1_phy_s1_q11_1.png}
\begin{qoptions}
\opt{A}{AB}
\opt{B}{$\overline{AB}$}
\opt{C}{$\bar AB+A\bar B$}
\opt{D}{$\overline{A+B}$}
\end{qoptions}
\end{bankq}
\begin{bankq}[id=SII13-P1-PHY-S2-Q02, source={Score-II - Test 13 (06 Feb 2024)}, paper=1, qno=2, type=numerical, subtopic={Diodes}, answer={3}, answer_src=key, difficulty=8, diff_parts={K3 S1 I2 C2 F1}, diff_note={concepts: standard problem in this chapter, diodes, apply the relevant results | steps: apply the relevant results; compute}]
Power lost in zener diode in the given circuit is $51.2\times10^{-x}$ w, find x :
\bfig{110.0}{SII13_p1_phy_s2_q02_1.png}
\end{bankq}
\begin{bankq}[id=SII17-P1-PHY-S1-Q18, source={Score-II - Test 17 (13 Feb 2024)}, paper=1, qno=18, type=single, subtopic={Logic gates}, answer={Bonus}, answer_src=solved, working={No answer key available (source truncated) and the figure/data needed to solve this question was not recoverable; banked as Bonus.}, difficulty=7, diff_parts={K3 S1 I2 C2 F0}, diff_note={concepts: standard problem in this chapter, logic gates, apply the relevant results | steps: apply the relevant results; compute}]
The logic circuit shown has the input waveforms `A' and `B' as shown. Pick out the \textbf{CORRECT} output waveform :-
\bfig{154.1}{SII17_p1_phy_s1_q18_1.png}
\begin{qoptions}
\opt{A}{\bfigi{153.8}{SII17_p1_phy_s1_q18_A.png}}
\opt{B}{\bfigi{154.3}{SII17_p1_phy_s1_q18_2.png}}
\opt{C}{\bfigi{154.3}{SII17_p1_phy_s1_q18_3.png}}
\opt{D}{\bfigi{154.3}{SII17_p1_phy_s1_q18_4.png}}
\end{qoptions}
\end{bankq}
\begin{bankq}[id=SII17-P1-PHY-S2-Q07, source={Score-II - Test 17 (13 Feb 2024)}, paper=1, qno=7, type=numerical, subtopic={Diodes}, answer={Bonus}, answer_src=solved, working={No answer key available (source truncated) and the figure/data needed to solve this question was not recoverable; banked as Bonus.}, difficulty=8, diff_parts={K3 S1 I2 C2 F1}, diff_note={concepts: standard problem in this chapter, diodes, apply the relevant results | steps: apply the relevant results; compute}]
In the given circuit, the diode has a forward resistance of 8 $\Omega$ and infinite backward resistance. The value of R is :-
\bfig{154.1}{SII17_p1_phy_s1_q07_1.png}
\end{bankq}
\begin{bankq}[id=SII19-P1-PHY-S1-Q11, source={Score-II - Test 19 (20 Feb 2024)}, paper=1, qno=11, type=single, subtopic={Diodes}, answer={A}, answer_src=key, difficulty=4, diff_parts={K2 S1 I1 C1 F0}, diff_note={concepts: standard problem in this chapter, diodes | steps: apply the relevant results; compute}]
In a circuit ideal diodes are given, find current drawn from battery?
\bfig{87.0}{SII19_p1_phy_s1_q11_1.png}
\begin{qoptions}
\opt{A}{2.00 ampere}
\opt{B}{2.31 ampere}
\opt{C}{1.33 ampere}
\opt{D}{1.71 ampere}
\end{qoptions}
\end{bankq}
\begin{bankq}[id=SII19-P1-PHY-S2-Q08, source={Score-II - Test 19 (20 Feb 2024)}, paper=1, qno=8, type=numerical, subtopic={Diodes}, answer={44}, answer_src=key, difficulty=8, diff_parts={K3 S1 I2 C2 F1}, diff_note={concepts: standard problem in this chapter, diodes, apply the relevant results | steps: apply the relevant results; compute}]
In the given circuit, the diode has a forward resistance of 8 $\Omega$ and infinite backward resistance. The value of $R$ is :-
\par\nobreak
\bfig{100.0}{SII19_p1_phy_s2_q08_1.png}
\end{bankq}
\begin{bankq}[id=SII22-P1-PHY-S1-Q12, source={Score-II - Test 22 (14 Mar 2024)}, paper=1, qno=12, type=single, subtopic={Logic gates}, answer={B}, answer_src=key, difficulty=4, diff_parts={K2 S1 I1 C1 F0}, diff_note={concepts: standard problem in this chapter, logic gates | steps: apply the relevant results; compute}]
Output for the following Boolean circuit is :
\bfig{126.7}{SII22_p1_phy_s1_q12_1.png}
\begin{qoptions}
\opt{A}{$\overline{AB}$}
\opt{B}{$\overline{AB}+\overline{AB}$}
\opt{C}{$\overline{A\overline{B}}+\overline{\overline{A}B}$}
\opt{D}{$A+B+\overline{AB}$}
\end{qoptions}
\end{bankq}
\begin{bankq}[id=SII22-P1-PHY-S2-Q06, source={Score-II - Test 22 (14 Mar 2024)}, paper=1, qno=6, type=numerical, subtopic={Diodes}, answer={5}, answer_src=key, difficulty=8, diff_parts={K3 S1 I2 C2 F1}, diff_note={concepts: standard problem in this chapter, diodes, apply the relevant results | steps: apply the relevant results; compute}]
A zener diode, having breakdown voltage equal to 15 V, is used in a voltage regulator circuit shown in figure. The current through the zener diode is (in mA) :
\bfig{80.8}{SII22_p1_phy_s2_q06_1.png}
\end{bankq}
\begin{bankq}[id=SII24-P1-PHY-S1-Q20, source={Score-II - Test 24 (17 Mar 2024)}, paper=1, qno=20, type=single, subtopic={Logic gates}, answer={C}, answer_src=key, difficulty=4, diff_parts={K2 S1 I1 C1 F0}, diff_note={concepts: standard problem in this chapter, logic gates | steps: apply the relevant results; compute}]
The truth table for digital circuit shown in figure is :-
\bfig{108.3}{SII24_p1_phy_s1_q20_1.png}
\begin{qoptions}
\opt{A}{\begin{tabular}[t]{|c|c|c|}\hline A&B&Y\\\hline 0&0&0\\0&1&0\\1&0&0\\1&1&1\\\hline\end{tabular}}
\opt{B}{\begin{tabular}[t]{|c|c|c|}\hline A&B&Y\\\hline 0&0&1\\0&1&0\\1&0&0\\1&1&0\\\hline\end{tabular}}
\opt{C}{\begin{tabular}[t]{|c|c|c|}\hline A&B&Y\\\hline 0&0&1\\0&1&1\\1&0&1\\1&1&0\\\hline\end{tabular}}
\opt{D}{\begin{tabular}[t]{|c|c|c|}\hline A&B&Y\\\hline 0&0&0\\0&1&1\\1&0&1\\1&1&0\\\hline\end{tabular}}
\end{qoptions}
\end{bankq}
\begin{bankq}[id=SII27-P1-PHY-S1-Q07, source={Score-II - Test 27 (23 Mar 2024)}, paper=1, qno=7, type=single, subtopic={p-n junction}, answer={D}, answer_src=key, difficulty=4, diff_parts={K2 S1 I1 C1 F0}, diff_note={concepts: standard problem in this chapter, p-n junction | steps: apply the relevant results; compute}]
The mobility of electrons and holes in a sample of insentric germanium at room temperature are 0.88 and 0.12 m$^2$/V-s. If electron and holes densities each are equal to $5.0\times10^{19}/\mathrm{m^3}$ then conductivity of germanium will be (in ohm meter) :-
\begin{qoptions}
\opt{A}{2}
\opt{B}{4}
\opt{C}{6}
\opt{D}{8}
\end{qoptions}
\end{bankq}
\begin{bankq}[id=SII30-P1-PHY-S1-Q08, source={Score-II - Test 30 (27 Mar 2024)}, paper=1, qno=8, type=single, subtopic={Diodes}, answer={D}, answer_src=key, difficulty=4, diff_parts={K2 S1 I1 C1 F0}, diff_note={concepts: standard problem in this chapter, diodes | steps: apply the relevant results; compute}]
Which of the following circuits is reverse - biased ?
\begin{qoptions}
\opt{A}{\bfigi{55.8}{SII30_p1_phy_s1_q08_2.png}}
\opt{B}{\bfigi{47.5}{SII30_p1_phy_s1_q08_1.png}}
\opt{C}{\bfigi{48.3}{SII30_p1_phy_s1_q08_3.png}}
\opt{D}{\bfigi{77.2}{SII30_p1_phy_s1_q08_4.png}}
\end{qoptions}
\end{bankq}
\begin{bankq}[id=SII30-P1-PHY-S1-Q09, source={Score-II - Test 30 (27 Mar 2024)}, paper=1, qno=9, type=single, subtopic={Logic gates}, answer={B}, answer_src=key, difficulty=4, diff_parts={K2 S1 I1 C1 F0}, diff_note={concepts: standard problem in this chapter, logic gates | steps: apply the relevant results; compute}]
The truth table of the given circuit diagram is:
\bfig{108.3}{SII30_p1_phy_s1_q09_1.png}
\begin{qoptions}
\opt{A}{\begin{tabular}[t]{|c|c|c|}\hline A&B&Y\\\hline 0&0&1\\0&1&0\\1&0&0\\1&1&1\\\hline\end{tabular}}
\opt{B}{\begin{tabular}[t]{|c|c|c|}\hline A&B&Y\\\hline 0&0&0\\0&1&1\\1&0&1\\1&1&0\\\hline\end{tabular}}
\opt{C}{\begin{tabular}[t]{|c|c|c|}\hline A&B&Y\\\hline 0&0&0\\0&1&0\\1&0&0\\1&1&1\\\hline\end{tabular}}
\opt{D}{\begin{tabular}[t]{|c|c|c|}\hline A&B&Y\\\hline 0&0&1\\0&1&1\\1&0&1\\1&1&0\\\hline\end{tabular}}
\end{qoptions}
\end{bankq}
\begin{bankq}[id=SII31-P1-PHY-S1-Q17, source={Score-II - Test 31 (29 Mar 2024)}, paper=1, qno=17, type=single, subtopic={Logic gates}, answer={A}, answer_src=key, difficulty=4, diff_parts={K2 S1 I1 C1 F0}, diff_note={concepts: standard problem in this chapter, logic gates | steps: apply the relevant results; compute}]
Here signals A, B, C are shown, then the logic gate is :-
\bfig{129.4}{SII31_p1_phy_s1_q17_1.png}
\bfig{116.3}{SII31_p1_phy_s1_q17_2.png}
\begin{qoptions}
\opt{A}{AND Gate}
\opt{B}{OR Gate}
\opt{C}{NOR Gate}
\opt{D}{NAND Gate}
\end{qoptions}
\end{bankq}
\begin{bankq}[id=SII31-P1-PHY-S2-Q05, source={Score-II - Test 31 (29 Mar 2024)}, paper=1, qno=5, type=numerical, subtopic={Diodes}, answer={142}, answer_src=key, difficulty=8, diff_parts={K3 S1 I2 C2 F1}, diff_note={concepts: standard problem in this chapter, diodes, apply the relevant results | steps: apply the relevant results; compute}]
In given circuit, power (in W) of zener diode for which bulb brights at rated value (Nearly) is :-
\bfig{115.9}{SII31_p1_phy_s2_q05_1.png}
\end{bankq}
\begin{bankq}[id=SII41-P1-PHY-S2-Q05, source={Score-II - Test 41 (19 Mar 2025)}, paper=1, qno=5, type=numerical, subtopic={Diodes}, answer={5}, answer_src=key, difficulty=8, diff_parts={K3 S1 I2 C2 F1}, diff_note={concepts: standard problem in this chapter, diodes, apply the relevant results | steps: apply the relevant results; compute}]
The value of power dissipated across the zener diode ($V_z=15$ V) connected in the circuit as shown in the figure is $x\times10^{-1}$ watt.
\bfig{137.3}{SII41_p1_phy_s2_q05_1.png}
The value of $x$, to the nearest integer, is \underline{\hspace{2em}}.
\end{bankq}
\begin{bankq}[id=SII45-P1-PHY-S1-Q07, source={Score-II - Test 45 (26 Mar 2025)}, paper=1, qno=7, type=single, subtopic={Diodes}, answer={D}, answer_src=key, difficulty=4, diff_parts={K2 S1 I1 C1 F0}, diff_note={concepts: standard problem in this chapter, diodes | steps: apply the relevant results; compute}]
In the circuit with ideal diodes as shown, current (in A) through battery is :-
\bfig{100.4}{SII45_p1_phy_s1_q07_1.png}
\begin{qoptions}
\opt{A}{3}
\opt{B}{5}
\opt{C}{6}
\opt{D}{4}
\end{qoptions}
\end{bankq}
\begin{bankq}[id=SII47-P1-PHY-S2-Q04, source={Score-II - Test 47 (30 Mar 2025)}, paper=1, qno=4, type=numerical, subtopic={Logic gates}, answer={1}, answer_src=key, difficulty=8, diff_parts={K3 S1 I2 C2 F1}, diff_note={concepts: standard problem in this chapter, logic gates, apply the relevant results | steps: apply the relevant results; compute}]
For the following circuit and given inputs A and B, find output between $t_5$ to $t_6$.
\bfig{430.5}{SII47_p1_phy_s2_q04_1.png}
\end{bankq}
\begin{bankq}[id=SII47-P1-PHY-S2-Q05, source={Score-II - Test 47 (30 Mar 2025)}, paper=1, qno=5, type=numerical, subtopic={Diodes}, answer={10}, answer_src=key, difficulty=8, diff_parts={K3 S1 I2 C2 F1}, diff_note={concepts: standard problem in this chapter, diodes, apply the relevant results | steps: apply the relevant results; compute}]
An ideal diode is connected in a circuit with resistance R = 50$\Omega$ and V = 10 volt as shown in figure, maximum value of output voltage (in V) in the positive cycle will be :
\bfig{272.7}{SII47_p1_phy_s2_q05_1.png}
\end{bankq}

\section*{Score Advanced papers}
\begin{bankq}[id=SA04-P2-PHY-S1-Q01, source={Score Advanced - Test 04 (04 Jul 2021)}, paper=2, qno=1, type=single, subtopic={Logic gates}, answer={C}, answer_src=solved, working={Y=(A+B)(A+C)}, difficulty=4, diff_parts={K1 S2 I1 C1 F0}, diff_note={concepts: combination of two OR gates and an AND gate | steps: upper OR gives A+B; lower OR gives A+C; AND of the two}]
The output of given logic circuit is :-\par
\bfig{153.5}{SA04_p2_phy_s1_q01_1.png}
\begin{qoptions}
\opt{A}{A $\cdot$ (B + C)}
\opt{B}{A$\cdot$ (B$\cdot$ C)}
\opt{C}{(A + B)$\cdot$ (A + C)}
\opt{D}{A + B + C}
\end{qoptions}
\end{bankq}
\begin{bankq}[id=SA04-P2-PHY-S1-Q11, source={Score Advanced - Test 04 (04 Jul 2021)}, paper=2, qno=11, type=single, subtopic={Diodes}, answer={B}, answer_src=solved, working={(10-0.3)/10=0.97 A and (10-6)/10=0.4 A, total 1.37 A, about 1.3 A}, difficulty=6, diff_parts={K2 S2 I2 C1 F0}, diff_note={concepts: diodes with a breakdown voltage, top diode conducts | steps: top diode conducts; lower diode is in reverse breakdown at 6 V; add the branch currents}]
Reverse break down voltage for silicon diode is 6V for a particular doping concentration and forward break down voltage is 0.3V then find approximate current through battery in the given circuit\par
\bfig{95.0}{SA04_p2_phy_s1_q11_1.png}
\begin{qoptions}
\opt{A}{2 Amp}
\opt{B}{1.3 Amp}
\opt{C}{1 Amp}
\opt{D}{$\dfrac{\pi}{\sqrt{2}}$ Amp}
\end{qoptions}
\end{bankq}
\begin{bankq}[id=SA04-P2-PHY-S1-Q16, source={Score Advanced - Test 04 (04 Jul 2021)}, paper=2, qno=16, type=single, subtopic={Transistors}, answer={C}, answer_src=solved, working={IC=0.625 mA, beta=24, IB=26 uA}, difficulty=5, diff_parts={K2 S2 I1 C1 F0}, diff_note={concepts: base current in a common emitter circuit, transistors | steps: IC=0.5/800; beta=alpha/(1-alpha)=24; IB=IC/beta}]
A transistor is connected in common emitter (CE) configuration. The collector supply is 8 V and the voltage drop across a resistor of 800 $\Omega$ in the collector circuit is 0.5 V. If the current gain factor $\alpha$ is 0.96, find the base current.
\begin{qoptions}
\opt{A}{48 $\mu$A}
\opt{B}{44 $\mu$A}
\opt{C}{26 $\mu$A}
\opt{D}{15 $\mu$A}
\end{qoptions}
\end{bankq}
\begin{bankq}[id=SA04-P2-PHY-S1-Q20, source={Score Advanced - Test 04 (04 Jul 2021)}, paper=2, qno=20, type=single, subtopic={Logic gates}, answer={C}, answer_src=solved, working={(A'+B)AB=AB so Y=(AB)' which is NAND}, difficulty=5, diff_parts={K2 S2 I1 C1 F0}, diff_note={concepts: logic circuit with NOT and OR and AND and NOT, logic gates | steps: Y'=A'+B; AND with A and B gives AB; inverted}]
Circuit given below is that of\par
\bfig{154.7}{SA04_p2_phy_s1_q20_1.png}
\begin{qoptions}
\opt{A}{AND gate}
\opt{B}{OR gate}
\opt{C}{NAND gate}
\opt{D}{NOR gate}
\end{qoptions}
\end{bankq}
\begin{bankq}[id=SA05-P1-PHY-S1-Q02, source={Score Advanced - Test 05 (15 Jul 2021)}, paper=1, qno=2, type=single, subtopic={p-n junction}, answer={C}, answer_src=key, difficulty=2, diff_parts={K1 S1 I0 C1 F0}, diff_note={concepts: field at an unbiased junction | steps: diffusion leaves the n side positive and the p side negative}]
In a P--N Junction diode not connected to any circuit
\begin{qoptions}
\opt{A}{potential is the same every where.}
\opt{B}{the P - type side is at a higher potential than the N - type side.}
\opt{C}{there is an electric field at the junction directed from the N - type side to the P-type side.}
\opt{D}{there is an electric field at the junction directed from the P - type side to the N-type side.}
\end{qoptions}
\end{bankq}
\begin{bankq}[id=SA05-P1-PHY-S1-Q06, source={Score Advanced - Test 05 (15 Jul 2021)}, paper=1, qno=6, type=single, subtopic={Logic gates}, answer={C}, answer_src=key, difficulty=3, diff_parts={K1 S1 I1 C1 F0}, diff_note={concepts: analogue and digital signals | steps: continuous signals are analogue; two-level signal is digital}]
The time variations of signals are given as in A, B and C. Point out the true statement from the following.\par
\bfig{91.1}{SA05_p1_phy_s1_q06_A.png}
\bfig{77.4}{SA05_p1_phy_s1_q06_B.png}
\bfig{75.7}{SA05_p1_phy_s1_q06_C.png}
\begin{qoptions}
\opt{A}{A, B and C are analogue signals}
\opt{B}{A and C digital, but B is analogue signal}
\opt{C}{A and C are analogue but B is digital signal}
\opt{D}{A, B and C are digital signal}
\end{qoptions}
\end{bankq}
\begin{bankq}[id=SA05-P1-PHY-S2-Q05, source={Score Advanced - Test 05 (15 Jul 2021)}, paper=1, qno=5, type=numerical, subtopic={Logic gates}, answer={2}, answer_src=key, difficulty=5, diff_parts={K2 S1 I1 C1 F1}, diff_note={concepts: sum of three gate outputs, logic gates | steps: evaluate each gate}]
In the following combination of logic gates, the outputs of A, B and C are x, y and z respectively then the value of x + y + z is\par
\bfig{105.5}{SA05_p1_phy_s2_q05_A.png}
\bfig{104.9}{SA05_p1_phy_s2_q05_B.png}
\bfig{103.6}{SA05_p1_phy_s2_q05_C.png}
\end{bankq}
\begin{bankq}[id=SA41-P1-PHY-S4-Q04, source={Score Advanced - Test 41}, paper=1, qno=4, type=numerical, subtopic={Logic gates}, answer={8}, answer_src=key, difficulty=8, diff_parts={K3 S1 I2 C2 F1}, diff_note={concepts: standard problem in this chapter, logic gates, apply the relevant results | steps: apply the relevant results; compute}]
A signal of 0.1 kW is transmitted in a cable. The attenuation of cable is $-$5 dB per km and cable length is 20 km. The power received at receiver is 10$^{-x}$ W. The value of x is \underline{\hspace{1cm}}.\par
[Gain in dB = 10 log$_{10}\left(\dfrac{\mathrm P_0}{\mathrm P_\mathrm i}\right)$]
\end{bankq}
\end{document}
